\documentclass[tikz,border=8pt]{standalone}
\usepackage{tikz}
\usepackage{amsmath}
\usepackage{amssymb}
\usepackage{pgfplots}
\pgfplotsset{compat=1.18}
\usepackage{ctex}
\usetikzlibrary{calc,positioning,arrows.meta,matrix,trees}

% LC 48 矩阵原地旋转 90度顺时针
% 4x4 矩阵，两步演化：先转置 → 再每行翻转

\begin{document}
\begin{tikzpicture}[
  every node/.style={font=\small},
  mcell/.style={draw=gray!60, rectangle, minimum size=0.8cm, thick, fill=white},
  hcell/.style={draw=blue!60, rectangle, minimum size=0.8cm, thick, fill=blue!15},
  gcell/.style={draw=green!60!black, rectangle, minimum size=0.8cm, thick, fill=green!15},
  rcell/.style={draw=orange!70!black, rectangle, minimum size=0.8cm, thick, fill=orange!12},
  ptr/.style={->,>=Stealth,thick},
]

\def\cs{0.92}

%% 标题
\node[font=\large\bfseries] at (8.0, 1.8)
  {矩阵原地旋转 90° 顺时针（LC 48）};
\node[font=\small,gray] at (8.0, 1.2)
  {两步：① 转置（沿主对角线对换）→ ② 每行翻转};

%% ─── 矩阵 A：原始（左侧）─────────────────────────────────
\node[font=\small\bfseries,blue!70] at (1.4, 0.5) {原始矩阵};

% [[ 1, 2, 3, 4],
%  [ 5, 6, 7, 8],
%  [ 9,10,11,12],
%  [13,14,15,16]]
% 对角线蓝色

\node[hcell,font=\scriptsize] at (0*\cs, 0)       {1};
\node[mcell,font=\scriptsize] at (1*\cs, 0)       {2};
\node[mcell,font=\scriptsize] at (2*\cs, 0)       {3};
\node[mcell,font=\scriptsize] at (3*\cs, 0)       {4};

\node[mcell,font=\scriptsize] at (0*\cs, -1*\cs)  {5};
\node[hcell,font=\scriptsize] at (1*\cs, -1*\cs)  {6};
\node[mcell,font=\scriptsize] at (2*\cs, -1*\cs)  {7};
\node[mcell,font=\scriptsize] at (3*\cs, -1*\cs)  {8};

\node[mcell,font=\scriptsize] at (0*\cs, -2*\cs)  {9};
\node[mcell,font=\scriptsize] at (1*\cs, -2*\cs)  {10};
\node[hcell,font=\scriptsize] at (2*\cs, -2*\cs)  {11};
\node[mcell,font=\scriptsize] at (3*\cs, -2*\cs)  {12};

\node[mcell,font=\scriptsize] at (0*\cs, -3*\cs)  {13};
\node[mcell,font=\scriptsize] at (1*\cs, -3*\cs)  {14};
\node[mcell,font=\scriptsize] at (2*\cs, -3*\cs)  {15};
\node[hcell,font=\scriptsize] at (3*\cs, -3*\cs)  {16};

% 图例标注（矩阵正下方）
\node[font=\tiny,blue!70] at (1.4, -3*\cs-0.55) {蓝色 = 主对角线元素};

%% 步骤1 箭头（水平，矩阵之间）
\draw[ptr,very thick,blue!70] (3.4, -1.5*\cs) -- (4.6, -1.5*\cs)
  node[midway,above,font=\scriptsize\bfseries,blue!70] {① 转置};

%% ─── 矩阵 B：转置后（中间）────────────────────────────────
\node[font=\small\bfseries,green!60!black] at (7.0, 0.5) {转置后};
\def\ox{5.1}

% 转置结果:
% [[ 1, 5, 9,13],
%  [ 2, 6,10,14],
%  [ 3, 7,11,15],
%  [ 4, 8,12,16]]

\node[hcell,font=\scriptsize] at (\ox+0*\cs, 0)       {1};
\node[gcell,font=\scriptsize] at (\ox+1*\cs, 0)       {5};
\node[gcell,font=\scriptsize] at (\ox+2*\cs, 0)       {9};
\node[gcell,font=\scriptsize] at (\ox+3*\cs, 0)       {13};

\node[gcell,font=\scriptsize] at (\ox+0*\cs, -1*\cs)  {2};
\node[hcell,font=\scriptsize] at (\ox+1*\cs, -1*\cs)  {6};
\node[gcell,font=\scriptsize] at (\ox+2*\cs, -1*\cs)  {10};
\node[gcell,font=\scriptsize] at (\ox+3*\cs, -1*\cs)  {14};

\node[gcell,font=\scriptsize] at (\ox+0*\cs, -2*\cs)  {3};
\node[gcell,font=\scriptsize] at (\ox+1*\cs, -2*\cs)  {7};
\node[hcell,font=\scriptsize] at (\ox+2*\cs, -2*\cs)  {11};
\node[gcell,font=\scriptsize] at (\ox+3*\cs, -2*\cs)  {15};

\node[gcell,font=\scriptsize] at (\ox+0*\cs, -3*\cs)  {4};
\node[gcell,font=\scriptsize] at (\ox+1*\cs, -3*\cs)  {8};
\node[gcell,font=\scriptsize] at (\ox+2*\cs, -3*\cs)  {12};
\node[hcell,font=\scriptsize] at (\ox+3*\cs, -3*\cs)  {16};

\node[font=\tiny,green!60!black] at (\ox+1.4, -3*\cs-0.55) {绿色 = 对换后的元素};

%% 步骤2 箭头
\draw[ptr,very thick,orange!70!black] (9.0, -1.5*\cs) -- (10.2, -1.5*\cs)
  node[midway,above,font=\scriptsize\bfseries,orange!70!black] {② 翻转};

%% ─── 矩阵 C：旋转结果（右侧）──────────────────────────────
\node[font=\small\bfseries,orange!70!black] at (12.7, 0.5) {旋转 90° 结果};
\def\px{10.7}

% 每行翻转结果:
% [[13, 9, 5, 1],
%  [14,10, 6, 2],
%  [15,11, 7, 3],
%  [16,12, 8, 4]]

\node[rcell,font=\scriptsize] at (\px+0*\cs, 0)       {13};
\node[rcell,font=\scriptsize] at (\px+1*\cs, 0)       {9};
\node[rcell,font=\scriptsize] at (\px+2*\cs, 0)       {5};
\node[rcell,font=\scriptsize] at (\px+3*\cs, 0)       {1};

\node[rcell,font=\scriptsize] at (\px+0*\cs, -1*\cs)  {14};
\node[rcell,font=\scriptsize] at (\px+1*\cs, -1*\cs)  {10};
\node[rcell,font=\scriptsize] at (\px+2*\cs, -1*\cs)  {6};
\node[rcell,font=\scriptsize] at (\px+3*\cs, -1*\cs)  {2};

\node[rcell,font=\scriptsize] at (\px+0*\cs, -2*\cs)  {15};
\node[rcell,font=\scriptsize] at (\px+1*\cs, -2*\cs)  {11};
\node[rcell,font=\scriptsize] at (\px+2*\cs, -2*\cs)  {7};
\node[rcell,font=\scriptsize] at (\px+3*\cs, -2*\cs)  {3};

\node[rcell,font=\scriptsize] at (\px+0*\cs, -3*\cs)  {16};
\node[rcell,font=\scriptsize] at (\px+1*\cs, -3*\cs)  {12};
\node[rcell,font=\scriptsize] at (\px+2*\cs, -3*\cs)  {8};
\node[rcell,font=\scriptsize] at (\px+3*\cs, -3*\cs)  {4};

\node[font=\tiny,orange!70!black] at (\px+1.4, -3*\cs-0.55) {橙色 = 最终结果};

%% ─── 验证说明（矩阵下方，统一区域）────────────────────────
\node[draw=gray!40,fill=white,thin,rounded corners=2pt,
      font=\scriptsize,inner sep=6pt,align=left]
  at (8.0, -4.2)
  {验证：原矩阵第0列（从下到上）= [13,9,5,1] = 旋转后第0行 \checkmark\quad
   原矩阵第3列（从下到上）= [16,12,8,4] = 旋转后第3行 \checkmark};

%% ─── 三个信息框（水平排列，y=-5.6）────────────────────
\node[draw=blue!40,fill=blue!5,thin,rounded corners=2pt,
      font=\small,inner sep=7pt,align=center]
  at (2.2, -5.6)
  {\textbf{步骤 ① 转置}\\[4pt]
   对 $0 \le i < j < n$\\[2pt]
   \texttt{swap(m[i][j], m[j][i])}};

\node[draw=orange!50,fill=orange!5,thin,rounded corners=2pt,
      font=\small,inner sep=7pt,align=center]
  at (8.0, -5.6)
  {\textbf{步骤 ② 翻转}\\[4pt]
   对每行执行\\[2pt]
   \texttt{row.reverse()}};

\node[draw=green!50,fill=green!5,thin,rounded corners=2pt,
      font=\small,inner sep=7pt,align=center]
  at (13.7, -5.6)
  {\textbf{等价变换}\\[4pt]
   顺时针 90°\\[2pt]
   $m'[j][n\!-\!1\!-\!i] = m[i][j]$};

%% 算法说明（y=-8.0）
\node[draw=gray!50,fill=white,rounded corners=3pt,font=\small,
      inner sep=8pt,align=left]
  at (8.0, -8.0)
  {\textbf{原地旋转 90° 顺时针 O($n^2$) 时间，O(1) 空间：}\\[4pt]
   先转置：对 $0 \le i < j < n$，交换 \texttt{matrix[i][j]} 与 \texttt{matrix[j][i]}\\[2pt]
   后翻转：对每行 \texttt{row}，执行 \texttt{row.reverse()}\\[2pt]
   逆时针 90°：先翻转每行，再转置（顺序相反）};

\end{tikzpicture}
\end{document}
